%% ============================================================
\section{Experiments}
\label{sec:experiments}
%% ============================================================

We evaluate DV-SLAM against three questions:
(1)~Can tightly-coupled LIO operate reliably with fewer than 500 points per
frame?
(2)~Is confidence-aware weighting essential, or is uniform ICP sufficient?
(3)~At what point count does sparse-regime LIO collapse, and does
confidence weighting extend the viable operating range?

%% ------------------------------------------------------------
\subsection{Experimental Setup}
\label{sec:setup}
%% ------------------------------------------------------------

\subsubsection{Datasets}

\textbf{ScanNet++~\cite{yeshwanth2023scannetpp} (primary).}
ScanNet++ (ICCV 2023) provides 460 high-fidelity indoor scenes captured on
an iPhone~13~Pro alongside FARO Focus~S70 laser-scan ground truth
(sub-millimetre accuracy).
Raw data include $256 \times 192$ 16-bit dToF depth, per-pixel confidence,
\texttt{CMDeviceMotion} IMU streams at 100\,Hz, RGB frames at
$1920 \times 1440$, and per-frame camera intrinsics---matching our pipeline's
input format exactly.
We evaluate on a held-out subset of sequences using the FARO mesh as
ground truth for map-quality scoring and the provided camera trajectory as
reference for ATE evaluation.

\textbf{Self-collected sequences (S1--S5, supplementary).}
We collect five closed-loop sequences on an iPhone~15~Pro~Max
(Table~\ref{tab:sequences}), logging raw IMU at 100\,Hz and dToF depth plus
per-pixel confidence at 30\,Hz.
COLMAP~\cite{schoenberger2016sfm} reconstructions from the co-captured RGB
stream serve as pseudo-ground truth for map quality; ARKit poses serve as
trajectory reference.

\begin{table}[t]
    \centering
    \caption{Self-collected evaluation sequences (iPhone~15~Pro~Max).
    COLMAP RGB pseudo-GT used for map quality; ARKit poses as trajectory
    reference.}
    \label{tab:sequences}
    \setlength{\tabcolsep}{4pt}
    \footnotesize
    \begin{tabular}{clcc}
        \toprule
        Seq. & Condition & Length & Purpose \\
        \midrule
        S1 & Normal indoor (room) & ${\sim}$20\,m & Baseline \\
        S2 & Long corridor       & ${\sim}$30\,m & Drift \\
        S3 & Fast rotation ($>$2\,rad/s) & ${\sim}$15\,m & IMU stress \\
        S4 & Glass / reflective surfaces & ${\sim}$15\,m & Confidence effect \\
        S5 & Bright outdoor & ${\sim}$20\,m & dToF saturation \\
        \bottomrule
    \end{tabular}
\end{table}

\subsubsection{Hardware}

All real-time experiments run on an iPhone~15~Pro~Max (Apple A17~Pro SoC,
single CPU thread, no GPU acceleration).
Table~\ref{tab:hardware} summarises the sensing and compute configuration.

\begin{table}[t]
    \centering
    \caption{Hardware configuration for all experiments.}
    \label{tab:hardware}
    \setlength{\tabcolsep}{4pt}
    \footnotesize
    \begin{tabular}{ll}
        \toprule
        Component & Specification \\
        \midrule
        Device & iPhone 15 Pro Max \\
        SoC    & Apple A17 Pro \\
        dToF sensor & LiDAR Scanner (rear-facing) \\
        dToF resolution & $256 \times 192$\,px, 30\,Hz \\
        Camera & 1920$\times$1440 RGB, 30\,Hz \\
        IMU    & 6-DoF \texttt{CMDeviceMotion}, 100\,Hz \\
        Processing & Single CPU thread; no GPU acceleration \\
        Dependencies & Eigen headers only (${\sim}$3\,MB) \\
        \bottomrule
    \end{tabular}
\end{table}

%% ------------------------------------------------------------
\subsection{Baselines}
\label{sec:baselines}
%% ------------------------------------------------------------

We compare against the following baselines:

\begin{enumerate}
    \item \textbf{ARKit VIO (platform baseline).}
    Apple's closed-source visual-inertial odometry pipeline, accessed via the
    \texttt{lowres\_wide.traj} file in ScanNet++ raw data.
    Requires continuous RGB input; provides the device-manufacturer upper
    bound on achievable accuracy.

    \item \textbf{Naive ICP (uniform weight, no IMU).}
    Point-to-plane ICP with uniform weights ($w_i^{\text{conf}} = 1$) and no
    IMU prediction.
    This baseline isolates the joint contribution of tight LIO coupling and
    confidence-aware weighting.

    \item \textbf{DV-SLAM (full, ours).}
    The complete pipeline: ESKF IMU prediction, Bundle-\&-Discard quality
    gate, confidence-weighted ICP update, and surface thinning on export.
\end{enumerate}

Existing dense-cloud LIO systems (FAST-LIO2~\cite{xu2022fastlio2},
DLIO~\cite{chen2023dlio}, LIO-SAM~\cite{shan2020liosam}) assume
$>$10K points per scan and are not directly comparable in the sparse dToF
regime ($<$500\,pts/frame).
For context, we note that na\"{i}ve application of FAST-LIO2 to these inputs
via a depth-to-\texttt{PointCloud2} bridge diverges immediately, confirming
the qualitative gap between the dense-cloud and sparse-dToF regimes.

%% ------------------------------------------------------------
\subsection{Metrics}
\label{sec:metrics}
%% ------------------------------------------------------------

\textbf{Trajectory.}
We report Absolute Trajectory Error (ATE) RMSE~\cite{sturm2012benchmark}
after Umeyama SE(3) alignment, computed with the \texttt{evo} toolbox.

\textbf{Map quality.}
We evaluate exported point clouds against ground-truth meshes using three
complementary metrics at a 5\,cm distance threshold:
\begin{itemize}
    \item \emph{Accuracy}: mean nearest-neighbor distance from estimated
    cloud to GT mesh (cm); measures reconstruction fidelity.
    \item \emph{Completeness}: mean nearest-neighbor distance from GT mesh
    to estimated cloud (cm); measures coverage.
    \item \emph{F-score}: harmonic mean of precision and recall at 5\,cm;
    balances both.
\end{itemize}

\textbf{Runtime.}
Per-stage wall-clock time on the A17~Pro (single CPU thread), reporting
mean and worst-case latency averaged over 100 consecutive keyframes.

%% ------------------------------------------------------------
\subsection{Main Results}
\label{sec:results}
%% ------------------------------------------------------------

Table~\ref{tab:main_comparison} reports results on ScanNet++ and the
self-collected sequences.
DV-SLAM (full) achieves the lowest ATE and highest F-score among camera-free
methods.
Naive ICP---lacking both IMU prediction and confidence weighting---shows
substantially higher drift, confirming that tight LIO coupling is necessary
for reliable operation in the sparse regime.
ARKit VIO, the camera-dependent platform baseline, serves as an upper-bound
reference.

\begin{table}[t]
    \centering
    \caption{Main comparison on ScanNet++ and self-collected S1--S5 (mean).
    ATE $=$ RMSE after SE(3) alignment via Umeyama.
    Acc./Comp.\ $=$ nearest-neighbor distances to GT mesh (cm).
    F-score at 5\,cm threshold.
    Camera column indicates whether RGB input is required.
    $\downarrow$: lower is better; $\uparrow$: higher is better.}
    \label{tab:main_comparison}
    \begin{tabular}{lcccc}
        \toprule
        Method & Cam.\ & ATE [cm] $\downarrow$ & Acc.\ [cm] $\downarrow$
               & F-score $\uparrow$ \\
        \midrule
        ARKit VIO (platform)               & \checkmark & -- & -- & -- \\
        Naive ICP ($w = 1$, no IMU)        & $\times$   & -- & -- & -- \\
        \textbf{DV-SLAM (ours)}            & $\times$   & -- & -- & -- \\
        \bottomrule
    \end{tabular}
\end{table}

%% ------------------------------------------------------------
\subsection{Ablation Study}
\label{sec:ablation}
%% ------------------------------------------------------------

Table~\ref{tab:ablation} isolates individual pipeline components on S1
(normal indoor, ${\sim}$20\,m closed loop).
Each row modifies exactly one aspect of the full system while holding all
other parameters fixed.

\begin{table}[t]
    \centering
    \caption{Ablation study on S1. Full $=$ Camera~+~Depth~+~IMU with all
    confidence-aware stages enabled.}
    \label{tab:ablation}
    \setlength{\tabcolsep}{4pt}
    \footnotesize
    \begin{tabular}{lccc}
        \toprule
        Configuration & ATE [cm] $\downarrow$ & Acc.\ [cm] $\downarrow$
                      & F-score $\uparrow$ \\
        \midrule
        Full (Camera + Depth + IMU)          & -- & -- & -- \\
        $-$ Camera (Depth + IMU only)        & -- & -- & -- \\
        $-$ Confidence ($w_i = 1$)           & -- & -- & -- \\
        $-$ B\&D (uniform voxel DS)          & -- & -- & -- \\
        $-$ TLS (L2 loss only)               & -- & -- & -- \\
        ${\sim}$100 pts/frame (stride 12)    & -- & -- & -- \\
        ${\sim}$3000 pts/frame (stride 1)    & -- & -- & -- \\
        \bottomrule
    \end{tabular}
\end{table}

\textbf{Row 2 ($-$Camera).}
Removing RGB input while retaining dToF and IMU tests the camera-free
operating mode central to DV-SLAM's privacy and robustness claims.

\textbf{Row 3 ($-$Confidence).}
Setting $w_i = 1$ directly tests Proposition~\ref{prop:efficiency}: ATE
degradation should be consistent with the efficiency ratio
$\eta \approx 1.54$.

\textbf{Row 4 ($-$B\&D).}
Replacing Bundle-\&-Discard quality gating with standard uniform voxel
downsampling at the same target count tests whether quality-aware point
selection is necessary or whether any reduction to ${\sim}$300 points
suffices.

\textbf{Row 5 ($-$TLS).}
Replacing the truncated least-squares robust cost with standard $L_2$
isolates the contribution of outlier rejection to map quality.

\textbf{Rows 6--7 (point-count boundary).}
Row~6 (stride~12, ${\sim}$100 pts/frame) probes the collapse boundary.
Row~7 (stride~1, ${\sim}$3000 pts/frame, no quality gate) verifies that
brute-force inclusion of all available points does not substitute for
confidence-aware selection.

%% ------------------------------------------------------------
\subsection{Point-Count Sweep}
\label{sec:sweep}
%% ------------------------------------------------------------

To characterize the collapse boundary empirically---the first such analysis
for sparse-dToF LIO---we vary the back-projection stride
$s \in \{1, 4, 8, 12, 16, 24\}$ on S1, yielding approximately
$\{3000, 500, 200, 100, 50\}$ points per frame.
Each configuration is evaluated with confidence weighting ON and OFF.

\begin{figure}[t]
    \centering
    \todo{generate figure: point-count sweep on S1}
    % \includegraphics[width=\columnwidth]{figures/point_count_sweep.pdf}
    \caption{%
    \textbf{First empirical characterization of the minimum viable
    observation count for LIO on sparse dToF.}
    ATE RMSE (log scale) vs.\ points per frame (log scale,
    $x \in \{50, 100, 200, 500, 1\text{K}, 3\text{K}\}$) on S1.
    Solid: DV-SLAM with confidence weighting.
    Dashed: uniform weights ($w_i = 1$).
    Vertical line: collapse boundary, defined as the point count below which
    median ATE exceeds 50\,cm.
    Confidence-aware processing shifts the collapse boundary downward,
    extending the viable operating range.%
    }
    \label{fig:sweep}
\end{figure}

As shown in Fig.~\ref{fig:sweep}, without confidence weighting ICP collapses
below ${\sim}$--\,pts/frame.
Confidence-aware processing extends stable tracking to ${\sim}$--\,pts/frame,
providing a meaningful safety margin for hardware-constrained deployments.

%% ------------------------------------------------------------
\subsection{Runtime Analysis}
\label{sec:runtime}
%% ------------------------------------------------------------

Table~\ref{tab:runtime} reports per-stage latency on the A17~Pro (single
CPU thread, no GPU).
The full pipeline completes well within the 100\,ms keyframe budget.

\begin{table}[t]
    \centering
    \caption{Per-stage runtime on Apple A17 Pro, single CPU thread.
    Averaged over 100 consecutive keyframes.}
    \label{tab:runtime}
    \setlength{\tabcolsep}{4pt}
    \footnotesize
    \begin{tabular}{lcc}
        \toprule
        Stage & Mean [ms] & Max [ms] \\
        \midrule
        IMU prediction (per IMU sample)         & -- & -- \\
        Bundle \& Discard (quality gate)        & -- & -- \\
        Visual tracking (optional RGB)          & -- & -- \\
        ICP + ESKF update                       & -- & -- \\
        \midrule
        \textbf{Total per keyframe}             & \textbf{--} & \textbf{--} \\
        Peak memory [MB]                        & \multicolumn{2}{c}{--} \\
        \bottomrule
    \end{tabular}
\end{table}
